\documentclass[10pt,a4paper]{amsart}
\usepackage{amsmath,amssymb,amsthm,graphicx}
\usepackage{verbatim}
\newenvironment{reference}{\comment}{\endcomment}
%\newenvironment{reference}{}{}
\newenvironment{slogan}{\comment}{\endcomment}
\newenvironment{history}{\comment}{\endcomment}
% For commutative diagrams we use Xy-pic
\usepackage[all]{xy}

% We use 2cell for 2-commutative diagrams.
\xyoption{2cell}
% Theorem environments.
%
\theoremstyle{plain}
\newtheorem{theorem}[subsection]{Theorem}
\newtheorem{proposition}[subsection]{Proposition}
\newtheorem{lemma}[subsection]{Lemma}

\theoremstyle{definition}
\newtheorem{definition}[subsection]{Definition}
\newtheorem{example}[subsection]{Example}
\newtheorem{exercise}[subsection]{Exercise}
\newtheorem{situation}[subsection]{Situation}

\theoremstyle{remark}
\newtheorem{remark}[subsection]{Remark}
\newtheorem{remarks}[subsection]{Remarks}

\numberwithin{equation}{subsection}

% Macros
%
\def\lim{\mathop{\mathrm{lim}}\nolimits}
\def\colim{\mathop{\mathrm{colim}}\nolimits}
\def\Spec{\mathop{\mathrm{Spec}}}
\def\Hom{\mathop{\mathrm{Hom}}\nolimits}
\def\Ext{\mathop{\mathrm{Ext}}\nolimits}
\def\SheafHom{\mathop{\mathcal{H}\!\mathit{om}}\nolimits}
\def\SheafExt{\mathop{\mathcal{E}\!\mathit{xt}}\nolimits}
\def\Sch{\mathit{Sch}}
\def\Mor{\mathop{\mathrm{Mor}}\nolimits}
\def\Ob{\mathop{\mathrm{Ob}}\nolimits}
\def\Sh{\mathop{\mathit{Sh}}\nolimits}
\def\NL{\mathop{N\!L}\nolimits}
\def\CH{\mathop{\mathrm{CH}}\nolimits}
\def\proetale{{pro\text{-}\acute{e}tale}}
\def\etale{{\acute{e}tale}}
\def\QCoh{\mathit{QCoh}}
\def\Ker{\mathop{\mathrm{Ker}}}
\def\Im{\mathop{\mathrm{Im}}}
\def\Coker{\mathop{\mathrm{Coker}}}
\def\Coim{\mathop{\mathrm{Coim}}}

% Boxtimes
%
\DeclareMathSymbol{\boxtimes}{\mathbin}{AMSa}{"02}

%
% Macros for moduli stacks/spaces
%
\def\QCohstack{\mathcal{QC}\!\mathit{oh}}
\def\Cohstack{\mathcal{C}\!\mathit{oh}}
\def\Spacesstack{\mathcal{S}\!\mathit{paces}}
\def\Quotfunctor{\mathrm{Quot}}
\def\Hilbfunctor{\mathrm{Hilb}}
\def\Curvesstack{\mathcal{C}\!\mathit{urves}}
\def\Polarizedstack{\mathcal{P}\!\mathit{olarized}}
\def\Complexesstack{\mathcal{C}\!\mathit{omplexes}}
% \Pic is the operator that assigns to X its picard group, usage \Pic(X)
% \Picardstack_{X/B} denotes the Picard stack of X over B
% \Picardfunctor_{X/B} denotes the Picard functor of X over B
\def\Pic{\mathop{\mathrm{Pic}}\nolimits}
\def\Picardstack{\mathcal{P}\!\mathit{ic}}
\def\Picardfunctor{\mathrm{Pic}}
\def\Deformationcategory{\mathcal{D}\!\mathit{ef}}
\usepackage{cleveref}
\crefname{originallemma}{original lemma}{original lemmas}
\Crefname{originallemma}{Original lemma}{Original lemmas}
\makeatletter
\newlabel{lemma-local-cohomology-adjoint}{{lemma-local-cohomology-adjoint}{source}{}{sourcepointer-1}{}}
\newlabel{lemma-local-cohomology-adjoint@cref}{{[originallemma][0][]lemma-local-cohomology-adjoint}{[1][0][]source}{}{}{}}
\makeatother
\crefname{originalreference}{original reference}{original references}
\Crefname{originalreference}{Original reference}{Original references}
\makeatletter
\newlabel{more-algebra-lemma-break-long-exact-sequence-modules}{{more-algebra-lemma-break-long-exact-sequence-modules}{source}{}{sourcepointer-2}{}}
\newlabel{more-algebra-lemma-break-long-exact-sequence-modules@cref}{{[originalreference][0][]more-algebra-lemma-break-long-exact-sequence-modules}{[1][0][]source}{}{}{}}
\makeatother
\crefname{originalreference}{original reference}{original references}
\Crefname{originalreference}{Original reference}{Original references}
\makeatletter
\newlabel{more-algebra-lemma-compute-Rlim-modules}{{more-algebra-lemma-compute-Rlim-modules}{source}{}{sourcepointer-3}{}}
\newlabel{more-algebra-lemma-compute-Rlim-modules@cref}{{[originalreference][0][]more-algebra-lemma-compute-Rlim-modules}{[1][0][]source}{}{}{}}
\makeatother
\crefname{originalreference}{original reference}{original references}
\Crefname{originalreference}{Original reference}{Original references}
\makeatletter
\newlabel{more-algebra-lemma-when-derived-completion-is-completion}{{more-algebra-lemma-when-derived-completion-is-completion}{source}{}{sourcepointer-4}{}}
\newlabel{more-algebra-lemma-when-derived-completion-is-completion@cref}{{[originalreference][0][]more-algebra-lemma-when-derived-completion-is-completion}{[1][0][]source}{}{}{}}
\makeatother
\title{The derived completion of local cohomology as a derived limit}
\author{The Stacks Project Authors}
\date{stacks-project:dualizing.tex:local-snapshot-20261012; GNU FDL 1.2 or later}
\begin{document}
\maketitle
\noindent\textit{Editorial note.} All delivered mathematical text is the author's own, copied from the pinned Stacks Project chapter source sources/stacks-dualizing/pinned-original-source.tex; no statement or proof step was rewritten, reordered or removed.\par\medskip
\noindent\textit{Editorial note.} Preamble provenance: the Stacks Project keeps one single shared preamble for every chapter; this card takes it directly from the staged sources/stacks-dualizing/preamble.tex as three verbatim spans (lines 8-12, 14-18 and 154-217), so no derived file was written and nothing was added or redefined. The wrapper class and the base packages are supplied by the assembler.\par\medskip
\noindent\textit{Editorial note.} Bibliography provenance: neither the counted statement nor its proof contains a citation (this lane rescans the whole span before writing the card), so the card carries no bibliography and no entry was reconstructed.\par\medskip
\noindent\textit{Editorial note.} Selection provenance: this card was opened under the file-based control contract (CONTROL.md), with the mechanical pre-card gate run on its statement and proof interval before the card was written.\par\medskip
\noindent\textit{Allowed standard prerequisites (editorial).}\begin{itemize}
\item The computation of local cohomology with support in a closed subscheme by the extended Cech complex. Invoked at the start of the counted proof (source lines 2344-2350). Provides the explicit complex whose derived completion is computed.
\item The description of the derived completion of that complex as the complex of ordinary completions. Invoked inside the counted proof (source lines 2351-2358). Rewrites the derived completion termwise.
\item The identification of a derived limit with the limit whenever the transition maps are surjective. Invoked inside the counted proof (source lines 2363-2364). Turns the termwise description into the asserted derived limit.
\item The breaking of a derived limit into a short exact sequence of the first derived limit and the limit. Invoked inside the counted proof (source lines 2365-2366). Produces the displayed short exact sequences of part (2).
\end{itemize}\medskip
\ifdefined\hypertarget\hypertarget{sourcepointer-1}{}\fi
\noindent\textit{Source reference.} Original lemma lemma-local-cohomology-adjoint: sources/stacks-dualizing/pinned-original-source.tex, source lines 1466--1466. Navigation pointer only: the counted text refers to the Stacks item lemma-local-cohomology-adjoint, which is an existing labeled result of this same chapter and is not part of this card. The pointer carries the target's own label as its number and its definition line as its span; its content is not reproduced here.\par\medskip
\ifdefined\hypertarget\hypertarget{sourcepointer-2}{}\fi
\noindent\textit{Source reference.} Original reference more-algebra-lemma-break-long-exact-sequence-modules: sources/stacks-dualizing/pinned-original-source.tex, source lines 2366--2366. Navigation pointer only: the counted text refers to the Stacks item more-algebra-lemma-break-long-exact-sequence-modules, which belongs to another chapter of the project and is not part of this card. The pointer carries the target's own label as its number because this chapter's staged snapshot does not contain the other chapter's tag table; its content is not reproduced here.\par\medskip
\ifdefined\hypertarget\hypertarget{sourcepointer-3}{}\fi
\noindent\textit{Source reference.} Original reference more-algebra-lemma-compute-Rlim-modules: sources/stacks-dualizing/pinned-original-source.tex, source lines 2364--2364. Navigation pointer only: the counted text refers to the Stacks item more-algebra-lemma-compute-Rlim-modules, which belongs to another chapter of the project and is not part of this card. The pointer carries the target's own label as its number because this chapter's staged snapshot does not contain the other chapter's tag table; its content is not reproduced here.\par\medskip
\ifdefined\hypertarget\hypertarget{sourcepointer-4}{}\fi
\noindent\textit{Source reference.} Original reference more-algebra-lemma-when-derived-completion-is-completion: sources/stacks-dualizing/pinned-original-source.tex, source lines 2352--2352. Navigation pointer only: the counted text refers to the Stacks item more-algebra-lemma-when-derived-completion-is-completion, which belongs to another chapter of the project and is not part of this card. The pointer carries the target's own label as its number because this chapter's staged snapshot does not contain the other chapter's tag table; its content is not reproduced here.\par\medskip
\begin{lemma}
\label{lemma-completion-local}
Let $I$ and $J$ be ideals in a Noetherian ring $A$. Let $M$ be a finite
$A$-module. Set $Z =V(J)$. Consider the derived $I$-adic completion
$R\Gamma_Z(M)^\wedge$ of local cohomology. Then
\begin{enumerate}
\item we have $R\Gamma_Z(M)^\wedge = R\lim R\Gamma_Z(M/I^nM)$, and
\item there are short exact sequences
$$
0 \to R^1\lim H^{i - 1}_Z(M/I^nM) \to H^i(R\Gamma_Z(M)^\wedge) \to
\lim H^i_Z(M/I^nM) \to 0
$$
\end{enumerate}
In particular $R\Gamma_Z(M)^\wedge$ has vanishing cohomology
in negative degrees.
\end{lemma}

\begin{proof}
Suppose that $J = (g_1, \ldots, g_m)$.
Then $R\Gamma_Z(M)$ is computed by the complex
$$
M \to \prod M_{g_{j_0}} \to \prod M_{g_{j_0}g_{j_1}} \to
\ldots \to M_{g_1g_2\ldots g_m}
$$
by Lemma \ref{lemma-local-cohomology-adjoint}.
By More on Algebra, Lemma
\ref{more-algebra-lemma-when-derived-completion-is-completion}
the derived $I$-adic completion of
this complex is given by the complex
$$
\lim M/I^nM \to \prod \lim (M/I^nM)_{g_{j_0}} \to
\ldots \to \lim (M/I^nM)_{g_1g_2\ldots g_m}
$$
of usual completions. Since $R\Gamma_Z(M/I^nM)$ is computed by
the complex $ M/I^nM \to \prod (M/I^nM)_{g_{j_0}} \to
\ldots \to (M/I^nM)_{g_1g_2\ldots g_m}$ and since the
transition maps between these complexes are surjective,
we conclude that (1) holds by
More on Algebra, Lemma \ref{more-algebra-lemma-compute-Rlim-modules}.
Part (2) then follows from More on Algebra, Lemma
\ref{more-algebra-lemma-break-long-exact-sequence-modules}.
\end{proof}

\end{document}
